\par\Needspace{6\baselineskip}\medskip
\begin{center}\hypertarget{native-168-article121}{}\pdfbookmark[2]{121.}{article-121}121.\end{center}
Whenever approximate values of the unknowns are already available, we can therefore \SourceLeaf{169}derive the true values from them by the method just explained with all the precision desired. First we calculate the values of $X,Y$ corresponding to those approximate values $a,b$. If these do not already vanish of their own accord, we repeat the calculation with another pair of values $a',b'$ that differ little from the first, and then with a third pair $a'',b''$, unless by chance $X$ and $Y$ vanished with the second pair. We then obtain the true values by the formulas of the preceding article, insofar as the assumption on which those formulas rest does not differ perceptibly from the truth. To judge this more accurately, we repeat the calculation of the values of $X,Y$ with those corrected values. If this shows that the equations $X=0$, $Y=0$ are still unsatisfied, the calculation will certainly yield much smaller values of $X,Y$ than the first three hypotheses did. The resulting orbital elements will accordingly be far more accurate than those corresponding to the first hypotheses. If we wish to go further, it is best to discard the hypothesis that produced the largest differences and combine the other two with the fourth. We thus form a fifth pair of values of $x,y$ following the preceding article. In the same way, if it seems worthwhile, we may proceed to a sixth hypothesis and so on, until the equations $X=0$, $Y=0$ are satisfied as accurately as the logarithmic and trigonometric tables permit. It will very rarely be necessary, however, to go beyond the fourth pair of values unless the initial hypotheses were too far from the truth.
